% TODO make hw stylesheet
\documentclass[10pt]{article}
\usepackage[letterpaper,top=1in,left=1in,right=1in]{geometry}
\usepackage[dvipsnames,svgnames]{xcolor}
\usepackage[shortlabels]{enumitem}
\usepackage[framemethod=TikZ]{mdframed}
\usepackage{amsmath,amssymb,amsthm}
\usepackage[colorlinks]{hyperref}
\usepackage{multicol}
\usepackage{tabularx}
\usepackage{bbm}

\hypersetup{
    colorlinks=true,
    linkcolor=blue,
    filecolor=magenta,
    urlcolor=magenta,
    pdftitle={MAT 324 HW}
}

\usepackage{mathtools}
\usepackage{thmtools}
\usepackage[textsize=footnotesize]{todonotes}
\usepackage{titling}
\usepackage{cancel}

\reversemarginpar
\mdfdefinestyle{mdbluebox}{roundcorner=10pt,innerbottommargin=9pt,
linecolor=blue,backgroundcolor=TealBlue!5,}
\declaretheoremstyle[headfont=\sffamily\bfseries\color{MidnightBlue},
mdframed={style=mdbluebox},]{thmbluebox}

\mdfdefinestyle{mdbloobox}{frametitlefont=\bfseries,innerbottommargin=8pt,
nobreak=true,backgroundcolor=TealBlue!5,linecolor=blue,}
\declaretheoremstyle[headfont=\bfseries\color{MidnightBlue},
mdframed={style=mdbloobox},headpunct={\\[3pt]},postheadspace=0pt,]{thmbloobox}

\mdfdefinestyle{mdredbox}{frametitlefont=\bfseries,innerbottommargin=8pt,
nobreak=true,backgroundcolor=Salmon!5,linecolor=RawSienna,}
\declaretheoremstyle[headfont=\bfseries\color{RawSienna},
mdframed={style=mdredbox},headpunct={\\[3pt]},postheadspace=0pt,]{thmredbox}

\mdfdefinestyle{mdgreenbox}{linecolor=ForestGreen,backgroundcolor=ForestGreen!5,
linewidth=2pt,rightline=false,leftline=true,topline=false,bottomline=false,}
\declaretheoremstyle[headfont=\bfseries\sffamily\color{ForestGreen!70!black},
mdframed={style=mdgreenbox},headpunct={ --- },]{thmgreenbox}

\mdfdefinestyle{mdorangebox}{linecolor=RedOrange,backgroundcolor=RedOrange!5,
linewidth=2pt,rightline=false,leftline=true,topline=false,bottomline=false,}
\declaretheoremstyle[headfont=\bfseries\sffamily\color{RedOrange!70!black},
mdframed={style=mdorangebox},headpunct={ --- },]{thmorangebox}

\mdfdefinestyle{mdblackbox}{linecolor=black,backgroundcolor=RedViolet!5!gray!5,
linewidth=3pt,rightline=false,leftline=true,topline=false,bottomline=false,}
\declaretheoremstyle[mdframed={style=mdblackbox}]{thmblackbox}

\declaretheoremstyle[spaceabove=6pt,spacebelow=6pt]{boldhead}
\declaretheoremstyle[spaceabove=6pt,spacebelow=6pt,headfont=\normalfont\itshape]{italichead}

\declaretheorem[style=boldhead,name=Problem]{problem}
\declaretheorem[style=boldhead,name=Problem,numbered=no]{problem*}
\declaretheorem[style=boldhead,name=Setup]{setup} % for config geo
\declaretheorem[style=boldhead,name=Example,sibling=problem]{example}
\declaretheorem[style=boldhead,name=$\bigstar$ Problem,sibling=problem]{niceprob}
\declaretheorem[style=boldhead,name=Theorem,sibling=problem]{theorem}
\declaretheorem[style=boldhead,name=Corollary,sibling=theorem]{corollary}
\declaretheorem[style=boldhead,name=Proposition,sibling=theorem]{prop}
\declaretheorem[style=boldhead,name=Lemma,numberwithin=problem]{lemma}
\declaretheorem[style=boldhead,name=Theorem,numbered=no]{theorem*}
\declaretheorem[style=boldhead,name=Lemma,numbered=no]{lemma*}
\declaretheorem[style=boldhead,name=Claim,numberwithin=problem]{claim}
\declaretheorem[style=boldhead,name=Claim,numbered=no]{claim*}
\declaretheorem[style=boldhead,name=Remark,numberwithin=problem]{remark}
\declaretheorem[style=boldhead,name=Remark,numbered=no]{remark*}
\declaretheorem[style=boldhead,name=Definition,sibling=theorem]{definition}
\declaretheorem[style=boldhead,name=Definition,numbered=no]{definition*}

\providecommand{\ol}{\overline}
\providecommand{\ceil}[1]{\left\lceil #1 \right\rceil}
\providecommand{\floor}[1]{\left\lfloor #1 \right\rfloor}
\newcommand{\dd}{\,\mathrm{d}}
\providecommand{\ul}{\underline}
\providecommand{\eps}{\varepsilon}
\providecommand{\half}{\frac{1}{2}}
\providecommand{\CC}{\mathbb C}
\providecommand{\FF}{\mathbb F}
\providecommand{\II}{\mathbb I}
\providecommand{\NN}{\mathbb N}
\providecommand{\QQ}{\mathbb Q}
\providecommand{\RR}{\mathbb R}
\providecommand{\ZZ}{\mathbb Z}
\providecommand{\ind}{\mathbbm 1}
\providecommand{\dg}{^\circ}
\providecommand{\ii}{\item}
\providecommand{\setdiff}{\smallsetminus}
\providecommand{\alert}[1]{{\sffamily\textbf{\textcolor{blue}{#1}}}}
\providecommand{\inv}{^{-1}}
\providecommand{\mb}[1]{\mathbf{#1}}
\newcommand{\what}{\widehat}

\newenvironment{soln}{\begin{proof}[Solution]}{\end{proof}}

\providecommand{\pow}{\operatorname{Pow}}
\providecommand{\vocab}[1]{{\textbf{\textcolor{ForestGreen}{#1}}}}
\providecommand{\lcm}{\operatorname{lcm}}
\providecommand{\abs}[1]{\left\lvert #1\right\rvert}
\providecommand{\smabs}[1]{\lvert #1\rvert}
\providecommand{\innerprod}[1]{\langle\!\langle #1\rangle\!\rangle}
\providecommand{\norm}[1]{\left\| #1\right\|}
\providecommand{\smnorm}[1]{\| #1\|}
\providecommand{\gr}{\operatorname{gr}}

\usepackage{mathrsfs}

% place title at top right
\setlength{\droptitle}{-8ex}
\pretitle{\begin{flushright}\Large\bfseries}
\posttitle{\par\end{flushright}}
\preauthor{\begin{flushright}}
\postauthor{\end{flushright}}
\predate{\begin{flushright}}
\postdate{\end{flushright}}

\title{MAT 324 HW11}
\author{\large Aditya Pahuja}
\date{\large Due 12/04}
\begin{document}
\maketitle

\begin{problem}
    \begin{enumerate}[(a)]
        \ii (Axler 11B \#2)
	Use the result of Exercise 12(a) in Section 11A to show that
	\begin{equation*}
	    1 + \frac 1{3^2} + \frac 1{5^2} + \frac 1{7^2} + \cdots = \frac{\pi^2}8
	\end{equation*}
	\ii Deduce the same identity from the identity in Axler's 11.32.
    \end{enumerate}
\end{problem}

\begin{soln}
    (a) Let $f$ be as in Exercise 11A \#12(a).
    Since $\abs{f}^2 = 1$ a.e., $\norm{f}_2^2 = 1$.
    By Parseval's identity,
    \begin{equation*}
	1 = \norm{f}_2^2 = \sum_{n\in\ZZ}\smabs{\what f(n)}^2
	= 2\sum_{n=1}^\infty \frac{4}{(2n-1)^2\pi^2}
	= \frac 8{\pi^2}
	\left(1 + \frac 1{3^2} + \frac 1{5^2} + \frac 1{7^2} + \cdots\right),
    \end{equation*}
    which implies what we want.
    \\

    (b) We can subtract the even terms away from Axler's 11.32:
    \begin{equation*}
	\sum_{n=1}^\infty \frac 1{(2n-1)^2}
	= \sum_{n=1}^\infty \frac 1{n^2}
	- \sum_{n=1}^\infty \frac 1{(2n)^2}
	= \left(1 - \frac 14\right)\sum_{n=1}^\infty \frac1{n^2}
	= \frac 34\cdot\frac{\pi^2}6 = \boxed{\frac{\pi^2}8}.
    \end{equation*}
    There is no cause for concern with the first equality
    because all the sums above converge absolutely.
\end{soln}

\begin{problem}
    Use Green's, Stokes', or residue theorem to show that
    \begin{equation*}
	\frac 1\pi\int_{-\infty}^\infty \frac{y}{x^2 + y^2}e^{-2\pi itx}\dd x
	= e^{-2\pi y\abs t}
    \end{equation*}
    for all $y\in\RR^+$ and $t\in\RR^*$.
\end{problem}

\begin{soln}
    Let $f\colon\CC\to\CC$ be
    \begin{equation*}
	f(z) = \frac 1\pi\cdot\frac{y}{z^2 + y^2} e^{-2\pi itz},
    \end{equation*}
    noting that $f$ has two singularities at $z = \pm iy$.
    Assume $t \le 0$, so that $e^{-2\pi itz}$ is decreasing with
    the imaginary part of $z$.
    Let $R > y$ and let $\gamma_R\subseteq\CC$ be the boundary of the region
    $\{z\in\CC : \abs{z}\le R\text{ and } \operatorname{Im}z \ge 0\}$,
    oriented positively. Since $\abs{f(z)} < \frac M{\abs{z}^2}$
    for some constant $M$ and sufficiently large $\abs z$
    (with nonnegative imaginary part),
    the integral over the semicircle becomes arbitrarily small
    as $R\to\infty$, so
    \begin{equation*}
	\int_{-\infty}^\infty f(x)\dd x =
	\lim_{R\to\infty} \oint_{\gamma_R} f(z)\dd z.
    \end{equation*}
    Since $iy$ is a simple pole of $f$,
    the residue of $f$ at $iy$ can be computed as
    \begin{equation*}
	\lim_{z\to iy} (z-iy)f(z)
	= \lim_{z\to iy}\frac 1\pi\cdot\frac{y(z-iy)}{(z+iy)(z-iy)} e^{-2\pi itz}
	= \lim_{z\to iy}\frac 1\pi \cdot\frac{y}{z+iy} e^{-2\pi itz}
	= \frac 1{2\pi i}e^{2\pi ty}
	= \frac 1{2\pi i}e^{-2\pi\abs ty}
    \end{equation*}
    where the last equality is because $t\le 0$.
    By the residue theorem,
    \begin{equation*}
	\oint_{\gamma_R} f(z)\dd z = 2\pi i\cdot \frac 1{2\pi i}e^{-2\pi \abs ty},
    \end{equation*}
    implying the result. Similarly, if $t\ge 0$,
    then we instead use residue theorem on $\gamma_R'$,
    the reflection of $\gamma_R$ over the real axis, in the exact same fashion.
\end{soln}

\begin{problem}
    \begin{enumerate}[(a)]
        \ii Let $(X,\mathcal S,\mu)$ be a measure space, $n\in\NN$,
	$p_1,\dots,p_n,r\in(0,\infty]$ with $1/p_1 + \cdots+ 1/p_n = 1/r$,
	and $f_1,\dots,f_n\colon X\to[0,\infty]$ be measurable functions. Show that
	\begin{equation*}
	    \norm{f_1\cdots f_n}_r \le \norm{f_1}_{p_1}\cdots\norm{f_n}_{p_n}.
	\end{equation*}
	\ii Let $p,q,r\in[1,\infty)$ with $1/p + 1/q - 1 = 1/r$
	and $f,g\colon\RR\to\RR$ be measurable functions.
	Show that for all $x\in\RR$,
	\begin{equation*}
	    \abs{(f*g)(x)}^r \le
	    \norm{f}_p^{r-p}\norm{g}_q^{r-q}
	    \int_\RR\abs{f(y)}^p\abs{g(x-y)}^q\dd\lambda(y).
	\end{equation*}
	\ii Let $p,q,r\in[1,\infty]$ with $1/p + 1/q - 1 = 1/r$ and
	$f,g\colon\RR\to\RR$ be measurable functions. Show that
	\begin{equation*}
	    \norm{f*g}_r \le \norm{f}_p\norm{g}_q.
	\end{equation*}
    \end{enumerate}
\end{problem}

\begin{soln}
    (a) We prove the inequality by induction. The base case $n=2$
    is the result of HW09 \#1(a). Suppose that the statement holds
    for $n - 1$ functions $f_i$ and numbers $p_i$,
    and let $x\in(0,\infty]$ be such that
    $1/{p_2} + 1/{p_3} + \cdots + 1/{p_n} = 1/x$. Since $1/p_1 + 1/x = 1/r$,
    the statement for $2$ functions and for $n-1$ functions together gives
    \begin{equation*}
	\norm{f_1\cdots f_n}_r \le \norm{f_1}_{p_1}\norm{f_2\cdots f_n}_{x}
	\le \norm{f_1}_{p_1}\cdot \norm{f_2}_{p_2}\cdots\norm{f_n}_{p_n}.
    \end{equation*}
    \\

    (b) Given $x\in\RR$, define $g_x\colon\RR\to\RR$ by
    $g_x(t) = g(x-t)$ for all $t\in\RR$. Then, taking both sides to the $1/r$ power,
    the desired inequality can be rewritten as
    \begin{equation*}
	\smnorm{fg_x}_1 \le \smnorm{\abs{f}^{1-p/r}}_{pr/(r-p)}\cdot
	\smnorm{\abs{g_x}^{1-q/r}}_{qr/(r-q)}\cdot
	\smnorm{\abs{f}^{p/r}\abs{g_x}^{q/r}}_{r},
    \end{equation*}
    which holds by part (a), noting that we are allowed to apply (a) since
    \begin{equation*}
        \left(\frac{pr}{r-p}\right)\inv + \left(\frac{qr}{r-q}\right)\inv + r\inv
	= \left(\frac 1p - \frac 1r\right) + \left(\frac 1q - \frac 1r\right)
	+ \frac 1r = \frac 1p + \frac 1q - \frac 1r \overset?= 1
    \end{equation*}
    by the given relation on $p$, $q$, and $r$.
    \\

    (c) Integrate the inequality proven in (b)
    and factor out the constants on the right-hand side to get
    \begin{equation*}
	\norm{f*g}_r^r \le
	\norm{f}_p^{r-p}\norm{g}_q^{r-q} \int_\RR\left(
	\int_\RR\abs{f(y)}^p\abs{g(x-y)}^q\dd\lambda(y)\right)\dd\lambda(x).
    \end{equation*}
    The double integral on the right-hand side evaluates to
    \begin{align*}
        \int_\RR\left(\int_\RR\abs{f(y)}^p
	\abs{g(x-y)}^q\dd\lambda(y)\right)\dd\lambda(x)
	&= \int_\RR\abs{f(y)^p)}\left(
	\int_\RR\abs{g(x-y)}^q\dd\lambda(x)\right)\dd\lambda(y)\\
	&= \int_\RR\abs{f(y)^p}\dd\lambda(y)\int_\RR\abs{g(x)}^q\dd\lambda(x)
	= \norm{f}_p^p\norm{g}_q^q,
    \end{align*}
    where the first equality is justified by Tonelli's theorem,
    as the integrand is nonnegative. Substituting this into the inequality gives
    \begin{equation*}
	\norm{f*g}_r^r \le \norm{f}_p^r\norm{g}_q^r,
    \end{equation*}
    which is equivalent to the desired result.
\end{soln}

\pagebreak

\begin{problem}
    \begin{enumerate}[(a)]
        \ii Suppose $f,\what f\in L^1(\RR;\CC)$.
	Show that $f,\what f\in L^2(\RR;\CC)$.
	\ii Suppose $f\in L^1(\RR;\CC)$ is differentiable
	with $f'\in L^1(\RR;\CC)$. Show that $\what f\in L^2(\RR;\CC)$.
	\ii Suppose $f\in L^1(\RR;\CC)$ is twice continuously differentiable
	with $f',f''\in L^1(\RR;\CC)$. Show that $\what f\in L^1(\RR;\CC)$.
    \end{enumerate}
\end{problem}

\begin{soln}
    (a) By the Axler's 11.49, $\what f,\widehat{\widehat{f\,}}\!\!
    \in L^\infty(\RR;\CC)$, and $\hat{\hat f}(x) = f(-x)$ for all $x\in\RR$,
    so $f\in L^\infty(\RR;\CC)$ too. This means
    \begin{equation*}
	\int_{\RR}\abs{f}^2\dd\lambda
	\le \norm{f}_\infty \int_\RR\abs{f}\dd\lambda
	= \norm{f}_\infty\norm{f}_1 < \infty,
    \end{equation*}
    hence $f\in L^2(\RR;\CC)$, and
    by the same logic $\what f\in L^2(\RR;\CC)$.
    \\

    (b) I think $f,f'\in L^1$ should imply $f$ is bounded? Assuming so,
    by H\"older's inequality,
    \begin{equation*}
	\norm{f}_2^2 = \norm{f^2}_1 \le \norm{f}_1\norm{f}_\infty < \infty,
    \end{equation*}
    so $f\in L^1\cap L^2$. By Plancherel's theorem,
    $\smnorm{\what f}_2 = \norm{f}_2 < \infty$, so $\what f\in L^2$.
    \\

    (c) By Axler's 11.54, $\smabs{\what f(t)} = 4\pi^2 t^2\smabs{\what{f''}(t)}$
    for all $t\in\RR$. By Axler's 11.49, $\what f$ and $\what{f''}$
    are both in $L^\infty$. Putting this together,
    \begin{equation*}
	\int_\RR\smabs{\what f}\dd\lambda
	= \int_{[-1,1]}\smabs{\what f}\dd\lambda
	+ \int_{\RR\setminus[-1,1]}\smabs{\what f}
	\le 2\smnorm{\what f}_\infty
	+ \int_{\RR\setminus[-1,1]}\frac{\smabs{\what{f''}}}{4\pi^2t^2} < \infty,
    \end{equation*}
    where the integral of $\smabs{\what{f''}}/4\pi^2t^2$ is finite
    because $\what{f''}$ is bounded; thus $\what f\in L^1$.
\end{soln}

\begin{problem}
    \begin{enumerate}[(a)]
        \ii Let $f\in L^2(\RR;\CC)$. Show that $\mathcal F^2f(x) = f(-x)$
	a.e. on $\RR$.
	\ii Suppose $f\in L^1(\RR;\CC)\cap L^2(\RR;\CC)$ and
	$g(x) = \int_\RR f(t) e^{2\pi ixt}\dd\lambda(t)$ for $x\in\RR$.
	Show that $g(x)\in\CC$ is well-defined for all $x\in\RR$,
	$g\in L^2(\RR;\CC)$, and $\mathcal Fg = f$.
    \end{enumerate}
\end{problem}

\begin{soln}
    (a) Suppose $f\in L^1\cap L^2$
    and, for $y>0$, let $P_y$ be the Poisson kernel.
    By \#3(c),
    \[
	\norm{f*P_y}_p\le \norm{f}_1\norm{P_y}_p < \infty
    \]
    for all $p\in[1,\infty]$ because $P_y\in L^p$
    for all $p\in[1,\infty]$, as shown in class,
    and $\what{f*P_y} = \what f*\what{P_y}\in L^p$
    because it is the product of an $L^\infty$ function with an $L^p$ function.
    In particular, the previous sentence shows that
    $f*P_y,\mathcal F(f*P_y)\in L^1\cap L^2$, so Fourier inversion implies that
    $\mathcal F^2(f*P_y)(x) = f*P_y(-x)$ a.e.
    Taking the limit in $L^2$ as $y\to 0^+$ and applying
    the result of \#7(a) implies that
    $\mathcal F^2 f(x) = f(-x)$ a.e.\ for all $f\in L^1\cap L^2$,
    which implies the result for all $f\in L^2$ because
    $L^1\cap L^2$ is dense in $L^2$ and $\mathcal F\colon L^2\to L^2$ is continuous.
    \\

    (b) The integral defining $g$ is well-defined because $f\in L^1$; explicitly,
    \begin{equation*}
	\abs{g(x)} \le \int_\RR\abs{f(t)e^{2\pi ixt}}\dd\lambda(t)
	= \norm{f}_1 < \infty.
    \end{equation*}
    Furthermore, $g(-x) = \what f(x) = \mathcal Ff(x)$ for all $x\in\RR$
    by definition of the Fourier transform,
    so Plancherel's theorem implies $\norm{g}_2 = \norm{f}_2 < \infty$.
    Also, using (a), we get that for all $x\in\RR$,
    \begin{equation*}
        f(x) = \mathcal F^2f(-x) = \mathcal Fg(x).\qedhere
    \end{equation*}
\end{soln}

\pagebreak
\begin{problem}
    \begin{enumerate}[(a)]
	\ii Let $f,g\in L^2(\RR;\CC)$. Show that $fg\in L^1(\RR;\CC)$ and
	$\what{fg} = (\mathcal Ff)*(\mathcal Fg)$.
	\ii Let $f,g\in L^2(\RR;\CC)$ with $f*g\in L^2(\RR;\CC)$.
	Show that $\mathcal F(f*g) = (\mathcal Ff)(\mathcal Fg)$ a.e. on $\RR$.
    \end{enumerate}
\end{problem}

\begin{soln}
    (a) Consider the convolution map $-*-\colon L^2\times L^2\to L^\infty$,
    which is well-defined because $\norm{f*g}_\infty \le \norm{f}_2\norm{g}_2$
    by \#3(c). For $f,g\in L^2$,
    each map $f*-\colon L^2\to L^\infty$, $-*g\colon L^2\to L^\infty$
    is linear (because convolution is bilinear) and bounded
    (because $\norm{f*g}_\infty \le \norm{f}_2\norm{g}_2$),
    so convolution is continuous in each component,
    whence it is easy to see that the convolution map
    on $L^2\times L^2$ (with the product metric
    $d((f_1,g_1),(f_2,g_2)) = \norm{f_1-f_2}_2 + \norm{g_1-g_2}_2$) is continuous.

    Since $\mathcal F\colon L^2\to L^2$, $\what{\bullet}\colon L^1\to L^\infty$,
    and the pointwise multilpication map $L^2\times L^2\to L^1$ are continuous
    (the second one by Riemann-Lebesgue and the latter by H\"older/Cauchy-Schwarz),
    it suffices to show that
    $(f,g)\mapsto (\mathcal Ff)*(\mathcal Fg)$
    agrees with $(f,g)\mapsto\what{fg}$ on the dense subset
    $(L^1\cap L^2)\times(L^1\cap L^2)\subseteq L^2\times L^2$. Indeed,
    \begin{equation*}
	\mathcal F(\mathcal Ff*\mathcal Fg)(-x)
	= (\mathcal F^2f\cdot\mathcal F^2g)(-x)
	= f(x)g(x)
	\iff 
	\mathcal F^2(\mathcal Ff*\mathcal Fg)(-x)
	= (\mathcal Ff*\mathcal Fg)(x)
	= \what{fg}(x),
    \end{equation*}
    for all $x\in\RR$,
    using the result of (b) for the leftmost equality
    (the proof of (b) below doesn't rely on (a)).

    (b) As shown in (a), convolution is continuous on $L^2\times L^2$;
    thus by density it suffices to verify the given identity for
    $f,g\in L^1\cap L^2$. In this case $\norm{f*g}_1 \le \norm{f}_1\norm{g}_1$
    by \#3(c), so $\what{f*g}$ is well-defined, meaning
    \begin{equation*}
	\mathcal F(f*g) = \what{f*g} = \what f\what g =
	(\mathcal Ff)(\mathcal Fg),
    \end{equation*}
    so we are done.
\end{soln}

\begin{problem}
    Let $f\in L^1(\RR;\CC)$.
    \begin{enumerate}[(a)]
        \ii For any $y > 0$, let $P_y$ be the Poisson kernel of Axler's 11.68.
	Show that $P_y*f, \what{P_y*f}\in L^p(\RR;\CC)$ for all $p\in[1,\infty]$,
	$P_y*f$ converges to $f$ as $y\to 0^+$
	with respect to $\norm{\bullet}_p$ for all $p\in[1,\infty)$
	such that $f\in L^p(\RR;\CC)$, and $\what{P_y*f}$ converges to $\what f$
	as $y\to 0^+$ with respect to $\norm{\bullet}_p$ for all $p\in[1,\infty]$
	such that $\what f\in L^p(\RR;\CC)$.
	\ii Suppose $\what f\in L^2(\RR;\CC)$. Show that $f\in L^2(\RR;\CC)$.
    \end{enumerate}
\end{problem}

\begin{soln}
    (a) Recall that $P_y, \what{P_y}\in L^p$ for all $p\in[1,\infty]$ and $y>0$
    (as shown in class).
    Using \#3(c),
    \[
	\norm{P_y*f}_p \le \norm{P_y}_p\norm{f}_1 < \infty,
    \]
    so $P_y*f\in L^p$. Furthermore $\what{P_y*f} = \what{P_y}\what f \in L^p$
    because it is the product of the $L^p$ function $\what{P_y}$
    with the bounded function $\what f$.

    Convergence of $P_y*f$ to $f$ when $f\in L^p$ for $p\in[1,\infty)$
    is exactly shown in Axler's 11.74.
    To show that $\what{P_y*f}$ converges to $\what f$ as $y\to 0^+$ if $p<\infty$,
    we compute
    \begin{equation*}
	\lim_{y\to 0^+} \smnorm{\what f - \what{P_y}\what f}_p^p
	= \lim_{y\to0^+}\int_\RR \smabs{\what f}^p\cdot\smabs{1-\what{P_y}}^p\dd\lambda
	= 0
    \end{equation*}
    with the last equality following from dominated convergence,
    where the integrand converges to and is uniformly bounded by
    $\smabs{\what f}^p$ as $y\to 0^+$.
    If $p=\infty$ then we note that $\what{P_y} = e^{-2\pi y\abs{t}}$
    converges uniformly to $1$ as $y\to0^+$, so, since $\what f\in L^\infty$
    is bounded almost everywhere, $\what{P_y}\what f$ converges
    with respect to the $L^\infty$ norm to $\what f$.
    \\

    (b) I reduced this to showing that
    $\norm{f*P_y}_2$ converges to $\norm{f}_2$ as $y\to 0^+$
    if $\norm{f}_2 = \infty$ (trying to show that
    $f\notin L^2$ implies $\what f\notin L^2$)
    but I don't see how to get the convergence
    (likely missing some silly bound on $f*P_y$ or something).
\end{soln}










\end{document}
